% === Begin Label Data ===
% ID: 854
% 难度星级: 
% 标签: 
% 备注: 
% 组卷引用次数: 0
% === End  Label Data ===

\begin{problem}{2016}{G}{四川卷（理）}{9}{导数}
设直线 $ l_1 $，$ l_2 $ 分别是函数 $ f(x)=\begin{cases} -\ln x, & 0<x<1, \\ \ln x, & x>1, \end{cases} $ 图像上点 $ P_1 $，$ P_2 $ 处的切线，$ l_1 $ 与 $ l_2 $ 垂直相交于点 $ P $，且 $ l_1 $，$ l_2 $ 分别与 $ y $ 轴相交于点 $ A $，$ B $，则 $ \triangle PAB $ 的面积的取值范围是 （\hspace{1cm}） (\hspace{1cm})
\begin{choices}
\choice{{$(0,1)$}}
\choice{{$(0,2)$}}
\choice{{$(0,+\infty)$}}
\choice{{$(1,+\infty)$}}
\end{choices}
\end{problem}

\begin{answer}
A
\end{answer}

\begin{solutions}
设 $ P_1(x_1,\ln x_1) $，$ P_2(x_2,-\ln x_2) $（不妨设 $ x_1>1 $，$ 0<x_2<1 $），则由导数的几何意义易得切线 $ l_1 $，$ l_2 $ 的斜率分别为 $ k_1=\dfrac{1}{x_1} $，$ k_2=-\dfrac{1}{x_2} $．

由已知得 $ k_1k_2=-1 $，$ \therefore x_1x_2=1 $，$ x_2=\dfrac{1}{x_1} $，切线 $ l_1 $ 的方程分别为 $ y-\ln x_1=\dfrac{1}{x_1}(x-x_1) $，切线 $ l_2 $ 的方程为 $ y+\ln x_2=-\dfrac{1}{x_2}(x-x_2) $，即 $ y-\ln x_1=-x_1\left(x-\dfrac{1}{x_1}\right) $．分别令 $ x=0 $ 得 $ A(0,-1+\ln x_1) $，$ B(0,1+\ln x_1) $．

又 $ l_1 $ 与 $ l_2 $ 的交点为 $ P\left(\dfrac{2x_1}{1+x_1^2},\ln x_1+\dfrac{1-x_1^2}{1+x_1^2}\right) $，

因为 $ x_1>1 $，所以 $ S_{\triangle PAB}=\dfrac{1}{2}\bigl|y_A-y_B\bigr|\cdot\bigl|x_P\bigr|=\dfrac{2x_1}{1+x_1^2}<\dfrac{1+x_1^2}{1+x_1^2}=1 $，$ \therefore 0<S_{\triangle PAB}<1 $，故选 A．
\end{solutions}